% GENERATED FILE. Edit canonical lessons, then run npm run generate:books.
% canonical-source-hash: fnv1a64:bbd94039a9d482b4
% canonical-lessons: SA-C467-sistah, SA-C467-upahasah, SA-C467-dhrstah, SA-C467-catuh, SA-C467-priyamvadah

\chapter{Good and Bad Manners}
\label{ch:sa-good-and-bad-manners}
\hlchaptermodality{467}

\begin{chapteropening}
\textbf{By the end of this chapter:} \emph{I can say whether someone is polite or rude.}

Complete the last lesson of chapter 467: I can say whether someone is polite or rude.
\end{chapteropening}

\section[\texorpdfstring{\sk{शिष्टः} (śiṣṭaḥ)}{शिष्टः (śiṣṭaḥ)}]{\sk{शिष्टः} (śiṣṭaḥ) — well-bred; cultured}

\label{lesson:SA-C467-sistah}



\begin{quote}
Before the new one: say the Sanskrit for childish, then the Sanskrit for understanding.
\end{quote}

\subsection*{You'll want to know: \sk{शिष्टः}}
\textbf{\sk{शिष्टः}} — \emph{śiṣṭaḥ} — \textquotedblleft{}well-bred; cultured\textquotedblright{}.

\begin{quote}
\sk{शिष्टः} \sk{बालः} \sk{अतिथिं} \sk{नमति।}
\end{quote}

\begin{quote}
\emph{śiṣṭaḥ bālaḥ atithiṃ namati.}
\end{quote}

\begin{quote}
The well-bred boy bows to the guest.
\end{quote}

\begin{cousinweb}
\textbf{\sk{शिष्टः}} (\emph{śiṣṭaḥ}) is built on the root √śās, \textquotedblleft{}to instruct, to train\textquotedblright{}: one who has been taught; MW gives \textquotedblleft{}disciplined, cultured, educated, wise\textquotedblright{} (MW).
\end{cousinweb}

\begin{grammarlens}[title={What you've built}]
One more word for describing what you lived and felt.
\end{grammarlens}

\subsection*{Your turn}
\begin{itemize}
\raggedright
  \item \emph{Say it:} \emph{śiṣṭaḥ}
  \item \emph{Say it:} \emph{śiṣṭaḥ}, once more
  \item \emph{Say it:} \emph{śiṣṭaḥ bālaḥ atithiṃ namati.} — The well-bred boy bows to the guest.
  \item \emph{Recall:} say the Sanskrit for childish, then the Sanskrit for understanding, then say \emph{śiṣṭaḥ} again
\end{itemize}

\begin{tcolorbox}[breakable,colback=teal!4,colframe=teal!35!black,title={Before you move on}]
What does \sk{शिष्टः} mean? (\textquotedblleft{}Well-bred; cultured\textquotedblright{}.) Say it once more.
\end{tcolorbox}
\section[\texorpdfstring{\sk{उपहासः} (upahāsaḥ)}{उपहासः (upahāsaḥ)}]{\sk{उपहासः} (upahāsaḥ) — mockery; derision}

\label{lesson:SA-C467-upahasah}



\begin{quote}
Before the new one: say the Sanskrit for understanding, then the Sanskrit for well-bred.
\end{quote}

\subsection*{You'll want to know: \sk{उपहासः}}
\textbf{\sk{उपहासः}} — \emph{upahāsaḥ} — \textquotedblleft{}mockery; derision\textquotedblright{}.

\begin{quote}
\sk{अतिथेः} \sk{उपहासः} \sk{न} \sk{उचितः।}
\end{quote}

\begin{quote}
\emph{atitheḥ upahāsaḥ na ucitaḥ.}
\end{quote}

\begin{quote}
Mocking a guest is not proper.
\end{quote}

A masculine noun.

\begin{cousinweb}
\textbf{\sk{उपहासः}} (\emph{upahāsaḥ}) is built from \emph{upa}, \textquotedblleft{}towards\textquotedblright{}, and the root √has, \textquotedblleft{}to laugh\textquotedblright{}, the root of \textbf{\sk{हसति}} and \textbf{\sk{उपहसति}}; MW gives \textquotedblleft{}laughter, derision, mockery, jeer\textquotedblright{} (MW).
\end{cousinweb}

\begin{grammarlens}[title={What you've built}]
One more word for describing what you lived and felt.
\end{grammarlens}

\subsection*{Your turn}
\begin{itemize}
\raggedright
  \item \emph{Say it:} \emph{upahāsaḥ}
  \item \emph{Say it:} \emph{upahāsaḥ}, once more
  \item \emph{Say it:} \emph{atitheḥ upahāsaḥ na ucitaḥ.} — Mocking a guest is not proper.
  \item \emph{Recall:} say the Sanskrit for understanding, then the Sanskrit for well-bred, then say \emph{upahāsaḥ} again
\end{itemize}

\begin{tcolorbox}[breakable,colback=teal!4,colframe=teal!35!black,title={Before you move on}]
What does \sk{उपहासः} mean? (\textquotedblleft{}Mockery; derision\textquotedblright{}.) Say it once more.
\end{tcolorbox}
\section[\texorpdfstring{\sk{धृष्टः} (dhṛṣṭaḥ)}{धृष्टः (dhṛṣṭaḥ)}]{\sk{धृष्टः} (dhṛṣṭaḥ) — impudent}

\label{lesson:SA-C467-dhrstah}



\begin{quote}
Before the new one: say the Sanskrit for well-bred, then the Sanskrit for mockery.
\end{quote}

\subsection*{You'll want to know: \sk{धृष्टः}}
\textbf{\sk{धृष्टः}} — \emph{dhṛṣṭaḥ} — \textquotedblleft{}impudent\textquotedblright{}.

\begin{quote}
\sk{धृष्टः} \sk{बालः} \sk{वृद्धम्} \sk{उपहसति।}
\end{quote}

\begin{quote}
\emph{dhṛṣṭaḥ bālaḥ vṛddham upahasati.}
\end{quote}

\begin{quote}
The impudent boy laughs at the old man.
\end{quote}

\begin{cousinweb}
\textbf{\sk{धृष्टः}} (\emph{dhṛṣṭaḥ}) comes from the root √dhṛṣ, \textquotedblleft{}to dare\textquotedblright{}; MW gives \textquotedblleft{}bold, daring, confident, audacious, impudent\textquotedblright{} (MW).
\end{cousinweb}

\begin{grammarlens}[title={What you've built}]
One more word for describing what you lived and felt.
\end{grammarlens}

\subsection*{Your turn}
\begin{itemize}
\raggedright
  \item \emph{Say it:} \emph{dhṛṣṭaḥ}
  \item \emph{Say it:} \emph{dhṛṣṭaḥ}, once more
  \item \emph{Say it:} \emph{dhṛṣṭaḥ bālaḥ vṛddham upahasati.} — The impudent boy laughs at the old man.
  \item \emph{Recall:} say the Sanskrit for well-bred, then the Sanskrit for mockery, then say \emph{dhṛṣṭaḥ} again
\end{itemize}

\begin{tcolorbox}[breakable,colback=teal!4,colframe=teal!35!black,title={Before you move on}]
What does \sk{धृष्टः} mean? (\textquotedblleft{}Impudent\textquotedblright{}.) Say it once more.
\end{tcolorbox}
\section[\texorpdfstring{\sk{चाटुः} (cāṭuḥ)}{चाटुः (cāṭuḥ)}]{\sk{चाटुः} (cāṭuḥ) — flattery; sweet talk}

\label{lesson:SA-C467-catuh}



\begin{quote}
Before the new one: say the Sanskrit for mockery, then the Sanskrit for impudent.
\end{quote}

\subsection*{You'll want to know: \sk{चाटुः}}
\textbf{\sk{चाटुः}} — \emph{cāṭuḥ} — \textquotedblleft{}flattery; sweet talk\textquotedblright{}.

\begin{quote}
\sk{धूर्तः} \sk{मधुरं} \sk{चाटुं} \sk{वदति।}
\end{quote}

\begin{quote}
\emph{dhūrtaḥ madhuraṃ cāṭuṃ vadati.}
\end{quote}

\begin{quote}
The rogue speaks sweet flattery.
\end{quote}

A masculine noun.

\begin{cousinweb}
\textbf{\sk{चाटुः}} (\emph{cāṭuḥ}) MW gives \textquotedblleft{}pleasing or graceful words or discourse, flattery\textquotedblright{}; MW gives no derivation, so its origin is uncertain (MW).
\end{cousinweb}

\begin{grammarlens}[title={What you've built}]
One more word for describing what you lived and felt.
\end{grammarlens}

\subsection*{Your turn}
\begin{itemize}
\raggedright
  \item \emph{Say it:} \emph{cāṭuḥ}
  \item \emph{Say it:} \emph{cāṭuḥ}, once more
  \item \emph{Say it:} \emph{dhūrtaḥ madhuraṃ cāṭuṃ vadati.} — The rogue speaks sweet flattery.
  \item \emph{Recall:} say the Sanskrit for mockery, then the Sanskrit for impudent, then say \emph{cāṭuḥ} again
\end{itemize}

\begin{tcolorbox}[breakable,colback=teal!4,colframe=teal!35!black,title={Before you move on}]
What does \sk{चाटुः} mean? (\textquotedblleft{}Flattery; sweet talk\textquotedblright{}.) Say it once more.
\end{tcolorbox}
\section[\texorpdfstring{\sk{प्रियंवदः} (priyaṃvadaḥ)}{प्रियंवदः (priyaṃvadaḥ)}]{\sk{प्रियंवदः} (priyaṃvadaḥ) — affable; kindly spoken}

\label{lesson:SA-C467-priyamvadah}



\begin{quote}
Before the new one: say the Sanskrit for impudent, then the Sanskrit for flattery.
\end{quote}

\subsection*{You'll want to know: \sk{प्रियंवदः}}
\textbf{\sk{प्रियंवदः}} — \emph{priyaṃvadaḥ} — \textquotedblleft{}affable; kindly spoken\textquotedblright{}.

\begin{quote}
\sk{प्रियंवदा} \sk{मीरा} \sk{अतिथिना} \sk{सह} \sk{मधुरं} \sk{वदति।}
\end{quote}

\begin{quote}
\emph{priyaṃvadā mīrā atithinā saha madhuraṃ vadati.}
\end{quote}

\begin{quote}
Affable Mira chats sweetly with the guest.
\end{quote}

\begin{cousinweb}
\textbf{\sk{प्रियंवदः}} (\emph{priyaṃvadaḥ}) is built from \textbf{\sk{प्रियः}}, \textquotedblleft{}dear, pleasant\textquotedblright{}, and \emph{vada}, \textquotedblleft{}speaking\textquotedblright{}, from the root √vad of \textbf{\sk{वदति}}: one who says pleasant things; MW gives \textquotedblleft{}speaking kindly, agreeable, affable\textquotedblright{} (MW).
\end{cousinweb}

\begin{grammarlens}[title={What you've built}]
One more word for describing what you lived and felt.
\end{grammarlens}

\subsection*{Your turn}
\begin{itemize}
\raggedright
  \item \emph{Say it:} \emph{priyaṃvadaḥ}
  \item \emph{Say it:} \emph{priyaṃvadaḥ}, once more
  \item \emph{Say it:} \emph{priyaṃvadā mīrā atithinā saha madhuraṃ vadati.} — Affable Mira chats sweetly with the guest.
  \item \emph{Recall:} say the Sanskrit for impudent, then the Sanskrit for flattery, then say \emph{priyaṃvadaḥ} again
\end{itemize}

\begin{tcolorbox}[breakable,colback=teal!4,colframe=teal!35!black,title={Before you move on}]
What does \sk{प्रियंवदः} mean? (\textquotedblleft{}Affable; kindly spoken\textquotedblright{}.) Say it once more.
\end{tcolorbox}
